% Source: https://artofproblemsolving.com/community/c3h2333352_problemo
% Stable identity: aops-topic-2333352-q01
Let the square have side $s$ and centre $O$, and use coordinates with $O$ at the origin and sides parallel to the axes. A circle of radius $1$ tangent to the two sides at a corner has its centre at distance $1$ from both sides, e.g. $P=\left(\dfrac s2-1,\dfrac s2-1\right)$. The four such centres form a square centred at $O$, and the circle of radius $2$, being at distance $1+2=3$ from all four of them, is centred at $O$ (as in the figure).

\begin{asy}
usepackage("amsmath");
size(4.5cm);
defaultpen(fontsize(10pt));
real s=2+3*sqrt(2); real h=s/2; real a=h-1;
draw((-h,-h)--(h,-h)--(h,h)--(-h,h)--cycle,linewidth(0.9));
draw(circle((0,0),2));
draw(circle((a,a),1)); draw(circle((-a,a),1)); draw(circle((a,-a),1)); draw(circle((-a,-a),1));
pair O=(0,0), P=(a,a);
draw(O--P,red); draw(P--(h,a),blue); draw(P--(a,h),blue);
dot(O); dot(P);
label("$O$",O,SW); label("$P$",P,NW,fontsize(8pt));
label("$1$",(a+0.5,a),S,blue+fontsize(8pt));
label("$OP=3$",(0.55*a,0.55*a),NW,red+fontsize(8pt));
\end{asy}

External tangency gives $OP=1+2=3$, and $OP=\sqrt2\left(\dfrac s2-1\right)$. Hence
\[\dfrac s2-1=\dfrac{3}{\sqrt2}=\dfrac{3\sqrt2}{2},\qquad s=2+3\sqrt2,\]
and the area is
\[s^2=4+12\sqrt2+18=22+12\sqrt2.\]
The answer is $\boxed{\textbf{(B) }22+12\sqrt2}$.
